% CLAUDE REPORT::BRIDGE-CLASS-0001::2026-11-13T10:00-08:00
\documentclass[11pt]{article}
\usepackage{amsmath,amssymb,amsthm,hyperref}
\newtheorem{theorem}{Theorem}
\newtheorem{corollary}[theorem]{Corollary}
\newtheorem{proposition}[theorem]{Proposition}
\title{Where the Worst-Case Derivation Programme Sits: R\_mix341 Is At Least Extended Resolution, and the Bridge B3 Is a Reformulation for Every System With a Known Lower Bound\\(BRIDGE-CLASS-0001)}
\author{Claude (scheduled routine; non-LSPV contributor, \texttt{research/agents/claude/})}
\date{13 November 2026}
\begin{document}\maketitle

\noindent\textbf{Status summary.}
Theorem~1 (R\_mix341 p-simulates Extended Resolution with $W\le L_{\rm ER}$): \textbf{PROVED/EXACT}.
Corollary~2 (a superpolynomial $\overline W_{R_{\rm mix341}}$ would be a superpolynomial Extended Frege lower bound): \textbf{PROVED/EXACT} given the \textbf{KNOWN THEOREM} that EF p-simulates ER.
Theorem~3 (classification of B3): \textbf{PROVED/EXACT}, elementary.
Proposition~4 (B3 for systems at least as strong as ER, restated as a reflection condition): \textbf{PROVED/EXACT} given Cook's PV/EF translation (\textbf{KNOWN THEOREM}).
Exhaustive DEF/extension gate check: \textbf{COMPUTATIONAL} (68 checks).
Superpolynomial $\overline W_R$, B3 for $R_{\rm mix341}$, and P versus NP: \textbf{OPEN}.

\section{Shared target}
For a fixed sound, refutationally complete calculus $R$, with the best derivation chosen for free:
\[
W_R(F)=\min_{\pi:F\vdash_R\bot}\operatorname{cost}_R(\pi),\qquad
\overline W_R(N)=\max_{F\in\mathrm{UNSAT},\ |\langle F\rangle|\le N}W_R(F).
\]
When cost is the number of DAG nodes, $W_R$ is the Cook--Reckhow proof length of $R$, so $\overline W_R$ is the worst-case length of the shortest proof. Choosing the best derivation for free is the standard convention in proof complexity, and it is a legitimate way to get lower bounds. The weakest sufficient bridge is
\[
\textbf{B3}_R:\qquad \mathrm{SAT}\in\mathrm P\ \Longrightarrow\ \overline W_R(N)\le N^{O(1)}.
\]
\textbf{LB}$_R$ denotes the statement $\overline W_R(N)=N^{\omega(1)}$. The programme is to prove $\textbf{LB}_R\wedge\textbf{B3}_R$, which together give $\mathrm{SAT}\notin\mathrm P$.

\section{Main obstacles and literature checked this run}
\textbf{Obstacle 1:} proving LB for a calculus strong enough to make B3 plausible.
\textbf{Obstacle 2:} proving B3 for a \emph{fixed} calculus chosen in advance.
\begin{itemize}
\item Cook--Reckhow (1979): some proof system is polynomially bounded if and only if NP $=$ coNP. This is a \textbf{KNOWN THEOREM}. It says nothing about which system, so it does not give B3 for a preselected $R$.
\item Krajíček--Pudlák (1989): if E $=$ NE, a p-optimal proof system exists. By Book's theorem, P $=$ NP implies that every tally NP set is in P, so E $=$ NE. Hence SAT $\in$ P gives \emph{some} polynomially bounded, p-optimal system. That system is not known to be p-equivalent to any fixed $R$, so this is a \textbf{failed transfer} to B3$_R$.
\item EF simulates every proof system $P$ once EF is given axioms stating that $P$ is sound; Cook's theory PV corresponds to EF. This is a \textbf{directly transferable theorem}; it is used in Proposition~4.
\item Pich--Santhanam, ``Why are proof complexity lower bounds hard?'' (FOCS 2019): under hardness assumptions, natural hard-tautology distributions block EF lower bounds, or else EF cannot efficiently prove circuit lower bounds. This is a \textbf{barrier that applies directly} to LB$_{R_{\rm mix341}}$, by Corollary~2.
\item Arteche--Khaniki--Pich--Santhanam, ``Towards P $\neq$ NP from Extended Frege lower bounds'' (ECCC TR23-199, arXiv:2312.08163): EF lower bounds for EF augmented with tautological ``witnessing'' axioms imply SAT $\notin$ P/poly. This is the closest known published form of the LB-to-separation bridge. It needs a \emph{modified} system, which supports the conclusion below that B3 has to be built into the choice of $R$.
\item Transfer domains: the knot/link handoff (2026-11-09) does not supply an object map to proof length. It is an \textbf{analogy only}. No post-quantum technique applies; this is a \textbf{failed transfer}.
\end{itemize}

\section{Results}
$R_{\rm mix341}$ is the calculus of LSPV-2/341 (\texttt{research/agents/lspv-2/2026-11-02\_1216\_LSPV-2\_MIXED\_TSEITIN\_MIN\_DAG341.tex}). Its rules are RES, AND3, XOR, LOCAL-CONVERT and DEF. Every derived node costs one unit and has at most two antecedents. DEF introduces a fresh variable defined as \emph{any} two-input Boolean gate, and all of its defining relations belong to that single paid node.

\begin{theorem}[$R_{\rm mix341}\ge_p$ ER, with constant 1]
For every unsatisfiable CNF $F$, $W_{R_{\rm mix341}}(F)\le L_{\rm ER}(F)$. Here $L_{\rm ER}$ is the number of lines in a shortest Extended Resolution refutation. The extension rule is $z\leftrightarrow(\ell_1\wedge\ell_2)$, and each extension counts as at least one line.
\end{theorem}
\begin{proof}
Fix an ER refutation and translate it line by line. Source clauses are free in both systems. A resolution line becomes one RES node. An extension $z\leftrightarrow(\ell_1\wedge\ell_2)$, with its three Tseitin clauses, is the two-input gate $g(a,b)=\ell_1\wedge\ell_2$. This gate is one of the 16 gates DEF allows, so it becomes one DEF node. That node carries all the defining clauses, and later resolution steps cite it. DEF is conservative, and freshness is preserved because the ER variable is fresh. The translated DAG uses at most as many paid nodes as the ER proof has lines, and it derives $\bot$.
\end{proof}
The script \texttt{research/agents/claude/bridge-class/def\_er\_check.py} checks the gate step exhaustively. For all 16 gates, the Tseitin clauses determine $z$ uniquely as $g(a,b)$ on all four inputs (64 checks). All four literal-sign variants of the ER extension appear among the 16 gates (4 checks). Result: 68/68 PASS, recorded in \texttt{def\_er\_check\_results.json}. The run takes under one second, so no chunking was needed.

\begin{corollary}[LB for $R_{\rm mix341}$ is at least as hard as an EF lower bound]
If $\overline W_{R_{\rm mix341}}(N)=N^{\omega(1)}$ on CNF inputs, then Extended Resolution, and therefore Extended Frege, needs superpolynomially many lines on some family of CNFs. That would be a new superpolynomial lower bound for Extended Frege, and no such bound is known even for Frege.
\end{corollary}
\begin{proof}
By Theorem~1, a large $W$ forces a large $L_{\rm ER}$. Size is at least the number of lines, and EF p-simulates ER (\textbf{KNOWN THEOREM}).
\end{proof}
\noindent\emph{Consequence for the corpus.} Every exponential $W$ lower bound previously proved for a \emph{restricted} variant (binary-premise Tseitin bounds, Res(PC$_r$) bounds via Pang 2026) is a lower bound for a subsystem that \emph{lacks} DEF. None of them can transfer to the full $R_{\rm mix341}$ unless it also proves an EF lower bound. This matches the corpus's own finding that $W_{R_{\rm mix341}}(G_k)=O(N_k)$ on Tseitin, and it explains it. The Haken pigeonhole control has the same explanation: Cook showed ER refutes the pigeonhole principle in polynomial size.

\begin{theorem}[Classification of B3]
Let $R$ be any sound, complete proof system.
\begin{enumerate}
\item[(i)] If P $\neq$ NP, then B3$_R$ holds vacuously.
\item[(ii)] If LB$_R$ is a theorem, then B3$_R\iff$ SAT $\notin$ P.
\item[(iii)] Consequently, for Resolution, Res($k$), Cutting Planes, Polynomial Calculus, bounded-depth Frege, and every restricted subsystem of $R_{\rm mix341}$ that the corpus has lower-bounded, B3 is \emph{exactly} P $\neq$ NP. In the world-ledger's terms it is a reformulation, not progress.
\end{enumerate}
\end{theorem}
\begin{proof}
(i) The antecedent of B3$_R$ is false. (ii) If LB$_R$ holds, B3$_R$ has a true antecedent only when SAT $\in$ P, and then its conclusion is false. So B3$_R$ is equivalent to $\neg(\mathrm{SAT}\in\mathrm P)$. Combined with (i), this gives the equivalence. (iii) Exponential lower bounds are known theorems for each system listed.
\end{proof}
\noindent So the two halves of the programme cannot be separated in difficulty in the hoped-for way. Where LB is provable, B3 carries the full difficulty of P $\neq$ NP. Where B3 might be easier (systems at least as strong as ER), LB is the Extended Frege lower-bound problem itself (Corollary~2).

\begin{proposition}[B3 for systems at least as strong as ER, restated as a reflection condition]
Let $R\ge_p$ ER, for example $R_{\rm mix341}$. Then B3$_R$ is equivalent to the following: if SAT $\in$ P, then some polynomial-time decider $A$ has polynomial-size $R$-proofs of the propositional soundness tautologies $\|\mathrm{Sound}_A\|_n$, which say that if $A$ rejects an input $x$ with $|x|=n$, then $x$ is unsatisfiable. A sufficient condition is that PV proves $\forall x\,\mathrm{Sound}_A(x)$ for one such $A$.
\end{proposition}
\begin{proof}
($\Leftarrow$) Given $F$ of size $n$: first derive $\|\mathrm{Sound}_A\|_n$. Then derive the polynomial-size EF verification of $A$'s rejecting computation on $F$, and resolve with $F$ to reach $\bot$. ER simulates this up to a polynomial factor, and so does $R$ by Theorem~1. ($\Rightarrow$) If $\overline W_R$ is polynomial, every true $\|\mathrm{Sound}_A\|_n$ has a short proof in $R$. The PV clause is Cook's translation (\textbf{KNOWN THEOREM}).
\end{proof}
\noindent This is the same premise as LSPV-1/337's finite-reflection B3 premise, now stated for the \emph{unmodified} $R_{\rm mix341}$. No additional reflection track is needed, because ER already contains the evaluation machinery. The status remains \textbf{OPEN}: P $=$ NP being true does not make it provable in PV.

\section*{CORPUS REVIEWED}
Read: AGENTS.md (all directives, including the 2026-11-09 knot/link and world-ledger notes) and the root \texttt{main.tex} ownership markers and $R_{\rm mix341}$ references (lines 3723--3906). Also read LSPV-2/341 (rule set), LSPV-1/337 (finite-reflection bridge), the LSPV-1/482 B1/B2 separation summary in the LSPV-5/464 outbox, and the LSPV-5/462 and /466 tables (Haken PHP, with ER as a mandatory upper-bound control). The most recent coordination handoffs were listed but not reread in full.

\section*{DUPLICATION / CONTRADICTION AUDIT}
Repository search found ER used as a \emph{control} (LSPV-5/462, /466) and Ext-PC$\surd$ noted as simulating EF (LSPV-5/465 outbox). It found no explicit statement that $R_{\rm mix341}$ itself p-simulates ER with constant 1, and no B3 classification theorem. No contradiction: Theorem~1 is consistent with every recorded upper bound on $R_{\rm mix341}$.

\section*{WHAT CHANGED RELATIVE TO ROOT MAIN}
Root \texttt{main.tex} has no Claude ownership region, so it was \textbf{not edited}. LSPV-5, as arbiter, should integrate two points: (a) LB$_{R_{\rm mix341}}$ is at least as hard as an EF lower bound; (b) B3 for systems with known lower bounds is a reformulation of P $\neq$ NP.

\section*{ADOPTED / REJECTED / RETIRED PRIOR COMPONENTS}
Adopted: $R_{\rm mix341}$ as the fixed calculus, and the LSPV-1/337 reflection premise. Rejected as a route: proving B3 for any system with a known lower bound, which is a reformulation. Nothing is retired.

\section*{STALE OR OVERSTATED CLAIMS FOUND}
None found, but this is a caution: any future claim of a superpolynomial $\overline W$ for the full $R_{\rm mix341}$, including DEF, must be audited as an EF lower bound.

\noindent\textbf{P versus NP remains OPEN.}
\end{document}
% END_OF_VERIFIED_REPORT::CLAUDE::BRIDGE-CLASS-0001
